\section{Censo de firmas del emisor residual}
\label{app:firmas-emision}

Este censo especifica las fibras finitas utilizadas en la sección~\ref{sec:extension}.
Cada condición inicial se identifica por
\[
 \nu=4(9i+j)+d,\qquad 0\leq i,j\leq8,\quad
 d=0,1,2,3\ \longleftrightarrow\ N,E,S,O.
\]
La división euclidiana recupera \(d=[\nu]_4\), \(j=[\lfloor\nu/4\rfloor]_9\)
e \(i=\lfloor\nu/36\rfloor\). Por tanto, \(0\leq\nu<324\) enumera
cada condición exactamente una vez, y las parejas de tales identificadores
enumeran las \(104\,976\) condiciones de \(\Omega_{\rm seed}\).

Sea \(f_+(\nu)\), respectivamente \(f_\times(\nu)\), la firma de seis
componentes obtenida por la recurrencia de emisión. En cada paso activo se lee
la casilla antes de avanzar; los vectores de dirección son
\[
 v_N=(-1,0),\quad v_E=(0,1),\quad v_S=(1,0),\quad v_O=(0,-1).
\]
Las coordenadas se reducen módulo nueve y el vector se invierte tras los
pasos \(27\) y \(54\). Los pasos de fase cero no desplazan ninguno de los
cursores. Estas reglas, junto con \eqref{eq:emisor-seis}, determinan
\(f_+\) y \(f_\times\) sin un catálogo de valores objetivo.

La tabla da todas sus imágenes. La multiplicidad \(m_\sigma\) es el número
de identificadores en la fibra \(f_\sigma^{-1}(a)\); \(\nu_{\min}\) es su
menor elemento. Cada palabra conserva los ceros iniciales. La fibra completa
queda determinada por la condición
\[
 f_\sigma^{-1}(a)=\{\nu\in\{0,\ldots,323\}:f_\sigma(\nu)=a\}.
\]
\begin{center}\small
\begin{tabular}{@{}lrr@{\qquad}lrr@{}}
\toprule
\(f_+\) & \(m_+\) & \(\nu_{\min}\) & \(f_\times\) & \(m_\times\) & \(\nu_{\min}\)\\
\midrule
\texttt{122564} & 18 & 17 & \texttt{000000} & 108 & 8 \\
\texttt{122898} & 18 & 24 & \texttt{177393} & 8 & 1 \\
\texttt{212645} & 18 & 5 & \texttt{177555} & 8 & 54 \\
\texttt{212988} & 18 & 0 & \texttt{177771} & 8 & 11 \\
\texttt{221456} & 18 & 29 & \texttt{339555} & 8 & 6 \\
\texttt{221889} & 18 & 12 & \texttt{339771} & 8 & 15 \\
\texttt{456221} & 18 & 28 & \texttt{339933} & 8 & 59 \\
\texttt{465898} & 18 & 21 & \texttt{393177} & 8 & 0 \\
\texttt{546988} & 18 & 9 & \texttt{393393} & 8 & 45 \\
\texttt{564122} & 18 & 16 & \texttt{393555} & 8 & 40 \\
\texttt{645212} & 18 & 4 & \texttt{555177} & 8 & 52 \\
\texttt{654889} & 18 & 33 & \texttt{555339} & 8 & 4 \\
\texttt{889221} & 18 & 13 & \texttt{555393} & 8 & 41 \\
\texttt{889654} & 18 & 32 & \texttt{555555} & 24 & 84 \\
\texttt{898122} & 18 & 25 & \texttt{555717} & 8 & 14 \\
\texttt{898465} & 18 & 20 & \texttt{555771} & 8 & 18 \\
\texttt{988212} & 18 & 1 & \texttt{555933} & 8 & 24 \\
\texttt{988546} & 18 & 8 & \texttt{717555} & 8 & 12 \\
 & &  & \texttt{717717} & 8 & 21 \\
 & &  & \texttt{717933} & 8 & 19 \\
 & &  & \texttt{771177} & 8 & 9 \\
 & &  & \texttt{771339} & 8 & 13 \\
 & &  & \texttt{771555} & 8 & 16 \\
 & &  & \texttt{933339} & 8 & 57 \\
 & &  & \texttt{933555} & 8 & 26 \\
 & &  & \texttt{933717} & 8 & 17 \\
\bottomrule
\end{tabular}
\end{center}

La evaluación de la recurrencia para los \(324\) identificadores produce
las dos particiones de la tabla: \(18\cdot18=324\) y
\(24\cdot8+24+108=324\). El producto de sus fibras da
\(432\) fibras de tamaño \(144\), \(18\) de tamaño \(432\) y \(18\)
de tamaño \(1944\), por la inyectividad de \eqref{eq:acoplamiento}.
Finalmente, aplicar \eqref{eq:psi-catalogo} a ese producto de imágenes
da \(243\) palabras distintas. El censo es finito y pertenece al emisor
residual; el refinamiento coinductivo posterior conserva su propia ley.

